\documentclass[11pt]{article}
\usepackage{amsmath, amssymb, amsthm, enumerate}
\usepackage[margin=1in, top=1.2in, bottom=1.2in]{geometry}
\usepackage{parskip}
\usepackage[colorlinks=true, linkcolor=blue, citecolor=blue, urlcolor=blue]{hyperref}
\usepackage{cleveref}
\theoremstyle{definition}
\newtheorem{definition}{Definition}[section]
\newtheorem{theorem}{Theorem}[section]
\newtheorem{lemma}{Lemma}[section]
\newtheorem{proposition}{Proposition}[section]
\newtheorem{corollary}{Corollary}[section]
\theoremstyle{remark}
\newtheorem{remark}{Remark}[section]
\newtheorem{example}{Example}[section]
\setlength{\parindent}{0pt}
\setlength{\parskip}{6pt}
\begin{document}

\begin{center}
{\Large\textbf{2.1 Measure Preliminaries-1}}\\
\vspace{0.3cm}
{\normalsize\today}
\end{center}
\vspace{0.5cm}


% --- Page 1 ---
\begin{definition}\label{def:algebra_of_sets}
Let $X$ be a non-empty set. An \underline{algebra of sets} of $X$ is a non-empty collection $\mathcal{A}$ of subsets of $X$ closed under finite union and complement.
\end{definition}

\begin{definition}\label{def:sigma_algebra}
A \underline{$\sigma$-algebra} on $X$ is a collection $\mathcal{M}$ of subsets of $X$ satisfying
\begin{enumerate}[(i)]
    \item $\emptyset \in \mathcal{M}$,
    \item If $E \in \mathcal{M}$, then $E^c = X \setminus E \in \mathcal{M}$,
    \item If $\{E_n\}_{n \in \mathbb{N}}$ is a sequence of sets in $\mathcal{M}$, then $\bigcup_{n=1}^\infty E_n \in \mathcal{M}$.
\end{enumerate}
Now $(X, \mathcal{M})$ is a \underline{measurable space} and elements of $\mathcal{M}$ are called \underline{measurable sets}.
\end{definition}

% --- Page 2 ---
\begin{enumerate}
    \item $X \in \mathcal{M}$

    \item If $E_1, \ldots, E_N \in \mathcal{M}$ then $\bigcup_{i=1}^N E_i \in \mathcal{M}$

    \item If $E_1, \ldots, E_N \in \mathcal{M}$ then $\bigcap_{i=1}^N E_i \in \mathcal{M}$

    \item If $\{E_i\}_{i \in \mathbb{N}}$ is a sequence of sets in $\mathcal{M}$, then $\bigcap_{i=1}^\infty E_i \in \mathcal{M}$

    \item If $E, F \in \mathcal{M}$ then $E \setminus F \in \mathcal{M}$

    \begin{proof}
    \begin{enumerate}
        \item $\emptyset \in \mathcal{M}$ and $X = (\emptyset)^c \in \mathcal{M}$.

        \item Let $\{E_i\}_{i \in \mathbb{N}}$ be a countable collection of measurable sets where $E_i = \emptyset$ for $i > N$. Hence,
        \[
        \bigcup_{i=1}^\infty E_i = E_1 \cup \cdots \cup E_N \in \mathcal{M}.
        \]

        \item $(\bigcup_{i=1}^\infty E_i)^c = \bigcap_{i=1}^\infty E_i^c \in \mathcal{M}$

        \item Apply 4 to the collection of sets in 2.

        \item $E \setminus F = E \cap F^c \in \mathcal{M}$ as finite intersection
    \end{enumerate}
    \end{proof}
\end{enumerate}

% --- Page 3 ---
\begin{proposition}\label{prop:intersection_sigma_algebra}
If $\{M_{\alpha} \mid \alpha \in J\}$ is a collection of $\sigma$-algebra on $X$. Then $\bigcap_{\alpha \in J} M_{\alpha}$ is a $\sigma$-algebra.
\end{proposition}

\begin{proof}
Observe $\emptyset \in M_{\alpha}$ for all $\alpha \in J$; this implies
$\emptyset \in \bigcap_{\alpha \in J} M_{\alpha}$. If $E \in \bigcap_{\alpha \in J} M_{\alpha}$, then $E \in M_{\alpha}$ for all $\alpha \in J$.
Hence $E^c \in M_{\alpha}$ for all $\alpha \in J$, this implies
$E^c \in \bigcap_{\alpha \in J} M_{\alpha}$. If $\{E_n\}_{n \in \mathbb{N}}$ is a sequence of sets in
$\bigcap_{\alpha \in J} M_{\alpha}$, hence $\bigcup_{n \in \mathbb{N}} E_n \in M_{\alpha}$ for each $\alpha \in J$.
Hence $\bigcup_{n \in \mathbb{N}} E_n \in \bigcap_{\alpha \in J} M_{\alpha}$.
\end{proof}

\begin{definition}\label{def:sigma_algebra_generated}
In particular, let $S \subseteq \mathcal{P}(X)$, there exists a unique smallest $\sigma$-algebra on $X$ that contains $S$. This is called $\sigma$-algebra generated by $S$.
\end{definition}

% --- Page 4 ---
\begin{remark}
We show uniqueness by constructing a collection of $\sigma$-algebras containing $S$ and then taking an intersection over the family of $\sigma$-algebras. By the above proposition, the resulting collection is the smallest $\sigma$-algebra containing $S$.
\end{remark}

\begin{definition}
Let $X$ be a topological space. The $\sigma$-algebra generated by the topology on $X$ is called \textbf{Borel $\sigma$-algebra} on $X$ and denoted by $\mathcal{B}_X$. The members of the $\sigma$-algebra are called \textbf{Borel sets}.
\end{definition}

\begin{remark}
Note: $F_\sigma$ is a countable union of closed sets and $G_\delta$ is a countable intersection of open sets. These are examples of Borel sets.
\end{remark}

% --- Page 5 ---
\begin{proposition}\label{prop:borel_generation}
$\mathcal{B}_{\mathbb{R}}$ is generated by each where $\mathcal{B}_{1\mathbb{R}}$ is a Borel $\sigma$-algebra:

\begin{enumerate}
    \item[(a)] open intervals $\mathcal{E}_1$
    \item[(b)] closed intervals $\mathcal{E}_2$
    \item[(c)] half-open intervals $\mathcal{E}_3$
    \item[(d)] open rays $\mathcal{E}_4$
    \item[(e)] closed rays $\mathcal{E}_5$
\end{enumerate}
\end{proposition}

\begin{proof}
Observe $\mathcal{E}_1$, $\mathcal{E}_2$, $\mathcal{E}_4$ and $\mathcal{E}_5$ are open or closed sets. For elements $\mathcal{E}_3$, these are $G_\delta$ sets, i.e. $(a,b] = \bigcap_{n=1}^\infty (a,b+\frac{1}{n})$. Hence all are Borel sets. If $\mathcal{M}(\mathcal{E}_j)$ is $\sigma$-algebra generated by $\mathcal{E}_j$, then $\mathcal{M}(\mathcal{E}_j) \subseteq \mathcal{B}_{\mathbb{R}}$ for $j = 1, 2, 3, 4, 5$.
\end{proof}

% --- Page 6 ---
As any open set in $\mathbb{R}$ can be expressed as a union of countable open intervals. Hence $\mathcal{B}_{\mathbb{R}} \subseteq \mathcal{M}(\mathcal{E}_j)$.

Now observe for any open interval $(a,b)$,
\[
(a,b) = \bigcap_{n=1}^\infty \left[a + \frac{1}{n}, b - \frac{1}{n}\right] = \bigcap_{n=1}^\infty \left[a, b + \frac{1}{n}\right] \cap \left[a - \frac{1}{n}, b\right]
\]
\[
= (a, \infty) \cap \bigcap_{n=1}^\infty \left(-\infty, b + \frac{1}{n}\right)
\]
\[
= \bigcap_{n=1}^\infty \left[a - \frac{1}{n}, \infty\right) \cap (-\infty, b)
\]

Hence $\mathcal{M}(\mathcal{E}_j) = \mathcal{B}_{\mathbb{R}}$ for $j = 1, 2, 3, 4, 5$.

\begin{remark}\label{rem:open_sets_intervals}
Every open set in $\mathbb{R}$ is a countable union of open intervals.
\end{remark}

Let $U$ be an open set in $\mathbb{R}$. If $x \in U$ then $x$ is rational or irrational. Let $x$ be rational then construct
\[
I_x = \bigcup_{x \in I} I
\]
where $I$ are open intervals containing $x$ and contained in $U$. If $x$ is not a rational for some $\varepsilon > 0$, $x \in (a - \varepsilon, a + \varepsilon) \subseteq U$, choose a rational $q \in (a - \varepsilon, a + \varepsilon)$ and observe

% --- Page 7 ---
---

Let $x \in I_q$. Hence for $x \in U$ there exists $q \in \mathbb{Q}$ such that $x \in I_q$. This implies $U \subseteq \bigcup_{q \in \mathbb{Q}} I_q$. Observe that $I_q \subseteq U$ for all $q \in \mathbb{Q}$ hence
\[
\bigcup_{q \in \mathbb{Q}} I_q = U \quad \text{where } I_q \text{ are open intervals.}
\]

\begin{definition}\label{def:ring}
A family of sets $\mathcal{R} \subseteq \mathcal{P}(X)$ is called a \emph{ring} if it is closed under finite union and difference (i.e. if $E_1, \dots, E_n \in \mathcal{R}$ then $\bigcup_{i=1}^n E_i \in \mathcal{R}$ and if $E, F \in \mathcal{R}$ then $E \setminus F \in \mathcal{R}$).
A ring that is closed under countable unions is called a $\sigma$-ring.
\end{definition}

\begin{proposition}\label{prop:sigma_ring_intersection}
If $\mathcal{R}$ is a $\sigma$-ring, then it is closed under countable intersection.
\end{proposition}

\begin{proof}
For $F \in \mathcal{R}$ observe $\emptyset = F \setminus F \in \mathcal{R}$. Let $\{E_i\}_{i=1}^\infty$ be a collection of sets belonging to $\mathcal{R}$. Then
\[
\bigcap_{i=1}^\infty E_i = \bigcap_{i=1}^\infty \left(F \setminus (F \setminus E_i)\right) = F \setminus \bigcup_{i=1}^\infty (F \setminus E_i) \in \mathcal{R},
\]
since $\mathcal{R}$ is a $\sigma$-ring.
\end{proof}

---

% --- Page 8 ---
To \(\mathcal{R}\). If any of them are pairwise disjoint
\(\bigwedge_{i=1}^n E_i = \emptyset \in \mathcal{R}\). Assume \(\{E_i\}_{i=1}^n\) are
not pairwise disjoint. As \(\mathcal{R}\) is a \(\sigma\)-ring, let
\(\bigcup_{i=1}^n E_i = E \in \mathcal{R}\). Let \(E_i' = E \setminus E_i \in \mathcal{R}\) for
\(i \in \mathbb{N}\). Now \(E' = \bigcup_{i=1}^n E_i' \in \mathcal{R}\). Observe
\(\bigwedge_{i=1}^n E_i = E \wedge E' \in \mathcal{R}\).

\[
\square
\]

% --- Page 9 ---
\begin{definition}\label{def:measure}
Let $X$ be a set equipped with a $\sigma$-algebra $\mathcal{M}$. We define a \textbf{countably additive measure} $m$ such that:
\begin{enumerate}[(i)]
    \item $m: \mathcal{M} \to [0, \infty]$ is a well-defined function.
    \item $m(\emptyset) = 0$.
    \item If $\{E_j\}_{j=1}^{\infty}$ is a sequence of disjoint sets, then
    \[
    m\left(\bigcup_{j=1}^{\infty} E_j\right) = \sum_{j=1}^{\infty} m(E_j).
    \]
\end{enumerate}
Observe that (iii) implies finite additivity by taking $E_j = \emptyset$ for $j > n$.
\end{definition}

We define $(X, \mathcal{M}, m)$ as a \textbf{measure space}.

\begin{definition}\label{def:finite_measure}
If $m(X) < +\infty$, we call $m$ a \textbf{finite measure}.
\end{definition}

\begin{definition}\label{def:sigma_finite_measure}
If $X = \bigcup_{j=1}^{\infty} E_j$ where $E_j \in \mathcal{M}$ and $m(E_j) < +\infty$ for all $j$, then $m$ is a \textbf{$\sigma$-finite measure}.
\end{definition}

\begin{definition}\label{def:sigma_finite_set}
More generally, $E$ is a \textbf{$\sigma$-finite set} if $E = \bigcup_{j=1}^{\infty} E_j$, $E_j \in \mathcal{M}$ and $m(E_j) < +\infty$ for all $j$.
\end{definition}

% --- Page 10 ---
\begin{definition}\label{def:semifinite}
If for each $E \in \mathcal{M}$ with $m(E) = \infty$ and there exists $F \subseteq E$ such that $0 < m(F) < \infty$, then $m$ is \emph{semifinite}.
\end{definition}

\begin{exercise}
Every $\sigma$-finite measure is semifinite.
\end{exercise}

\begin{proof}
If $m(E) = \infty$, then observe $E = E \cap X = \bigcup_{j=1}^{\infty} (E \cap E_j)$.
This implies $E \cap E_j \subseteq E_j \subseteq X$ for some $j$. We know $m(E_j) < +\infty$, hence
$m(E \cap E_j) < m(E_j) < +\infty$. Hence we have found a subset of $E$, specifically $E \cap E_j$, which has measure of $E_j \cap E$ to be finite.
\end{proof}

% --- Page 11 ---
\begin{example}\label{ex:dirac_measure}
Let $X$ be any set and $x_0 \in X$ be a fixed point and $\mathcal{M} = \mathcal{P}(X)$. Let $m : \mathcal{M} \to \mathbb{R}$

\[
m(E) = \begin{cases}
1 & \text{if } x_0 \in E \\
0 & \text{if } x_0 \notin E
\end{cases}
\]

This is called the \textbf{Dirac measure}.
\end{example}

\begin{example}\label{ex:counting_measure}
Let $X$ be any set and $\mathcal{M} = \mathcal{P}(X)$. Let $m : \mathcal{M} \to [0, \infty]$

\[
m(E) = \begin{cases}
|E| & \text{if } |E| < +\infty \\
\infty & \text{otherwise}
\end{cases}
\]

where $| \cdot |$ is the cardinality function and $m$ is called the \textbf{counting measure}.
\end{example}

\begin{example}\label{ex:infinite_measure}
Let $X$ be any set and $\mathcal{M} = \mathcal{P}(X)$. Let $m : \mathcal{M} \to [0, \infty]$

\[
m(E) = +\infty \quad \text{for all } E \subset X.
\]
\end{example}

% --- Page 12 ---
\begin{lemma}
Suppose $X$ be a set and $\mathcal{M}$ is a $\sigma$-algebra equipped on $X$. Let $m: \mathcal{M} \to [0, \infty]$ be a measure, then

\begin{enumerate}[(i)]
    \item \textbf{[Monotonicity]}: If $A, B \in \mathcal{M}$ such that $A \subseteq B$, then
    \[
    m(A) \leq m(B).
    \]

    \item If $\{A_n\}_{n \in \mathbb{N}} \subseteq \mathcal{M}$, then
    \[
    m\left(\bigcup_{n \in \mathbb{N}} A_n\right) \leq \sum_{n \in \mathbb{N}} m(A_n).
    \]
\end{enumerate}
\end{lemma}

\begin{proof}
\begin{enumerate}[(i)]
    \item If $B \supseteq A$, then
    \[
    m(B) = m(A \cup (B \setminus A)) = m(A) + m(B \setminus A) \geq m(A).
    \]

    \item Let $E_1 = A_1$, and $E_k = A_k \setminus \bigcup_{i=1}^{k-1} A_i$ for $k \geq 1$. Then observe all $E_k$ are disjoint to each other, hence
    \[
    m\left(\bigcup_{k \in \mathbb{N}} E_k\right) = \sum_{k \in \mathbb{N}} m(E_k).
    \]
    Note that
    \[
    \bigcup_{n \in \mathbb{N}} A_n = \bigcup_{k \in \mathbb{N}} E_k \quad \text{and} \quad m(E_k) \leq m(A_k) \quad \text{for all } k \in \mathbb{N}.
    \]
    Hence we get
    \[
    m\left(\bigcup_{n \in \mathbb{N}} A_n\right) \leq \sum_{n \in \mathbb{N}} m(A_n).
    \]
\end{enumerate}
\end{proof}

% --- Page 13 ---
\begin{lemma}
Let $(X, \mathcal{M}, \mu)$ be a measure space.

\begin{enumerate}
    \item \textbf{Continuity from below:} If $\exists\, E_j \uparrow^\infty_j = 1 \in \mathcal{M}$ and $E_1 \subseteq E_2 \subseteq \cdots$, then
    \[
    \mu\left(\bigcup_{j=1}^\infty E_j\right) = \lim_{j \to \infty} \mu(E_j).
    \]

    \item \textbf{Continuity from above:} If $\exists\, E_j \downarrow^\infty_j = 1 \in \mathcal{M}$ and $E_1 \supseteq E_2 \supseteq \cdots$, then
    \[
    \mu\left(\bigcap_{j=1}^\infty E_j\right) = \lim_{j \to \infty} \mu(E_j),
    \]
    and $\mu(E_1) < +\infty$.

\end{enumerate}

\begin{proof}

    \begin{enumerate}
        \item[(a)] Set $E_0 = \emptyset$. Observe $\bigcup_{j=1}^\infty E_j = \bigcup_{j=1}^\infty E_j \setminus E_{j-1}$. Hence,
        \[
        \mu\left(\bigcup_{j=1}^\infty E_j\right) = \mu\left(\bigcup_{j=1}^\infty (E_j \setminus E_{j-1})\right)
        \]
        \[
        = \sum_{j=1}^\infty \mu(E_j \setminus E_{j-1}) = \lim_{n \to \infty} \sum_{j=1}^n \mu(E_j \setminus E_{j-1})
        \]
        \[
        = \lim_{n \to \infty} \mu(E_n).
        \]

        \item[(b)] Let $F_j = E_1 \setminus E_j$. Observe $F_1 \subseteq F_2 \subseteq \cdots$ and $\bigcup_{j=1}^\infty F_j = E_1 \setminus \bigcap_{j=1}^\infty E_j$. Then,
        \[
        \mu(E_1) = \mu\left(\bigcap_{j=1}^\infty E_j\right) + \lim_{n \to \infty} \mu(F_n).
        \]
    \end{proof}
\end{lemma}

% --- Page 14 ---
---
\[
\mu(E_i) = \mu\left(\bigcap_{i=1}^\infty E_j\right) + \lim_{n \to \infty} \mu(E_i \setminus E_j)
\]

\[
= \mu\left(\bigcap_{i=1}^\infty E_j\right) + \mu(E_i) - \lim_{n \to \infty} \mu(E_j).
\]

This implies
\[
\mu\left(\bigcap_{i=1}^\infty E_j\right) = \lim_{n \to \infty} \mu(E_j),
\]
as $\mu(E_i) < +\infty$, we can subtract it from both sides.

\begin{proof}
$\square$
\end{proof}

\begin{definition}\label{def:null_set}
If $(X, \mathcal{M}, \mu)$ is a measure space, a set $E \in \mathcal{M}$ such that $\mu(E) = 0$ is called a \underline{null set}.

If a statement is true about point $x \in X \setminus E$ for some null set $E$, then we say the statement holds \underline{almost everywhere} (a.e). (More precisely, we have $\mu$-null set and $\mu$-a.e.)
\end{definition}

\begin{corollary}\label{cor:null_set_union}
Any countable union of null sets is a null set, and subsets of null sets belonging in $\mathcal{M}$ are null sets.
\end{corollary}
---

% --- Page 15 ---
\begin{definition}\label{def:complete_measure_domain}
If a measure whose domain includes all subsets of a null set then the domain is called \textbf{complete}.
\end{definition}

\begin{theorem}\label{thm:measure_extension}
Suppose that $(X, \mathcal{M}, \mu)$ is a measure space. Let
\[
\mathcal{N} = \{N \in \mathcal{M} : \mu(N) = 0\}
\]
and
\[
\overline{\mathcal{M}} = \{E \cup F \mid E \in \mathcal{M} \text{ and } F \subseteq N \text{ for some } N \in \mathcal{N}\}.
\]
Then $\overline{\mathcal{M}}$ is a $\sigma$-algebra and there is a unique extension $\overline{\mu}$ of $\mu$ to complete measure $\overline{\mathcal{M}}$.
\end{theorem}

\begin{proof}
\textbf{(a) $\overline{\mathcal{M}}$ is a $\sigma$-algebra}

Since $\mathcal{M}$ and $\mathcal{N}$ are closed under countable unions, so is $\overline{\mathcal{M}}$. If $E \in \mathcal{M}$ and $F \subseteq N \in \mathcal{N}$ then $E \cup F \in \overline{\mathcal{M}}$. Now assume $E \cap N = \emptyset$ otherwise replace $F$ by $F \setminus E$ and $N$ by $N \setminus E$. Then
\[
E \cup F = (E \cup N) \cap (N^c \cup F),
\]
so
\[
(E \cup F)^c = (E \cup N)^c \cup (N^c \cup F)^c
\]
and observe
\end{proof}

% --- Page 16 ---
\[
(EUN)^c \in \mathcal{M} \text{ and } N \setminus F^c = N \setminus F \subseteq N \text{ and } N \setminus F \in \mathcal{N}.
\text{ Hence } (EUF)^c \in \overline{\mathcal{M}}.
\]

\begin{proof}
(cb) Define $\overline{\mu} : \overline{\mathcal{M}} \to [0, \infty]$, such that
\[
\overline{\mu}(EUF) = \mu(CE).
\]
We want to show well definedness. Let $E_1UF_1 = E_2UF_2$ where
$E_1, E_2 \in \mathcal{M}$ and $F_1, F_2 \in \mathcal{N}$. Now,
\[
E_1 \subseteq E_1UF_1 = E_2UF_2 \subseteq E_2UN_2.
\]
This implies
\[
\mu(CE_1) \leq \mu(E_2) + \mu(N_2) = \mu(CE_2).
\]
If $E_2 \subseteq E_1UN$, and this implies $\mu(E_2) \leq \mu(E_1)$. Hence
\[
\mu(CE_1) = \mu(CE_2).
\]
\end{proof}

(c) Uniqueness.

Let $m_1, m_2$ be extensions of $\mu$ as per the above extension. Now, $EUF \in \overline{\mathcal{M}}$. Observe
\[
m_1(EUF) = \mu(CE) = m_2(EUF')
\]
where $F'$ is any subset of a null set.

% --- Page 17 ---
\begin{definition}\label{def:outer_measure}
An outer measure on a non-empty set $X$ is a function $\mu^* \colon \mathcal{P}(X) \to [0, \infty]$ that satisfies,
\begin{enumerate}
    \item $\mu^*(\emptyset) = 0$.
    \item $\mu^*(A) \leq \mu^*(B)$ if $A \subseteq B$.
    \item $\mu^*\left(\bigcup_{n=1}^\infty A_n\right) \leq \sum_{n=1}^\infty \mu^*(A_n)$.
\end{enumerate}
\end{definition}

\begin{proposition}\label{prop:outer_measure_construction}
Let $\mathcal{E} \subseteq \mathcal{P}(X)$ and $\rho \colon \mathcal{E} \to \mathcal{P}(X)$ and $\rho \colon \mathcal{E} \to [0, \infty]$ be such that $\emptyset \in \mathcal{E}$, $X \in \mathcal{E}$ and $\rho(\emptyset) = 0$. For any $A \subset X$ we define
\[
\mu^*(A) = \inf\left\{\sum_{i=1}^\infty \rho(E_i) \mid E_i \in \mathcal{E} \text{ and } A \subseteq \bigcup_{i=1}^\infty E_i\right\}.
\]
Then $\mu^*$ is an outer measure.
\end{proposition}

% --- Page 18 ---
\begin{proof}
(a) Observe a countable collection of empty sets covers $A = \emptyset$ hence
\[
\mu^* \emptyset = \inf_{\substack{\mathcal{E} \\ \bigcup \mathcal{E} \supseteq \emptyset}} \sum_{E_n \in \mathcal{E}} \rho(E_n) \leq \sum_{n \in \mathbb{N}} \rho(\emptyset) = 0.
\]
Also $\rho(E_n) \geq 0$ for any $E_n \in \mathcal{E}$ hence
\[
\sum_{E_n \in \mathcal{E}} \rho(E_n) \geq 0 \quad \text{for any collection of intervals}
\]
therefore
\[
\inf_{\substack{\mathcal{E} \\ \bigcup \mathcal{E} \supseteq \emptyset}} \sum_{E_n \in \mathcal{E}} \rho(E_n) \geq 0.
\]
So we get $\mu^* \emptyset = 0$.

For any non-empty set $A \subseteq X$ there exists $\{E_j\}_{j=1}^\infty \subseteq \mathcal{E}$ such that $A \subseteq \bigcup_{j=1}^\infty E_j$ (we can set $E_j = X$ for all $j$).
\end{proof}

% --- Page 19 ---
\begin{proof}
(cb) If $A \subseteq B$ then if $\{E_j\}$ is a collection of sets covering $B$ then $\mathcal{B}$ covers $A$ as well.

If $\mathcal{B}$ is a cover of $B$, we construct
\[
\mathcal{A} = \{E_n \in \mathcal{B} \mid B_n \cap A \neq \emptyset\}
\]
and observe $\mathcal{A} \subseteq \mathcal{B}$. Hence
\[
\sum_{A_n \in \mathcal{A}} P(A_n) \leq \sum_{B_n \in \mathcal{B}} P(B_n).
\]
Therefore
\[
\mu^* A \leq \sum_{B_n \in \mathcal{B}} P(B_n).
\]
Observe $\mu^* A \leq \mu^* B$; otherwise $\mu^* A > \mu^* B$ then there exist a countable collection $\{E_j\} = E$ covering $B$ and also $A$ whose sum, i.e. $\sum P(B_n)$, is less than $\mu^* A$, which contradicts the fact that $\mu^* A$ is the infimum of sums of countable collections of sets in $E$ covering $A$.
\end{proof}

% --- Page 20 ---
\begin{proof}
(c) Suppose $\{A_j\}_{j=1}^\infty \subset \mathcal{P}(X)$ and fix $\delta > 0$.
For each $j \in \mathbb{N}$ we get $\{E_{j,k}\}_{k=1}^\infty \subset \mathcal{E}$ such that
\[
A_j \subset \bigcup_{k=1}^\infty E_{j,k} \quad \text{and} \quad \sum_{k=1}^\infty \rho(E_{j,k}) \leq \mu^*(A_j) + \frac{\delta}{2^j}.
\]
But if $A = \bigcup_{j=1}^\infty A_j$, we have $A \subset \bigcup_{j,k=1}^\infty E_{j,k}$ and
\[
\sum_{j,k} \rho(E_{j,k}) \leq \sum_j \rho(E_{j,k}) + \delta,
\]
whence
\[
\mu^*(A) \leq \sum_j \mu^*(A_j) + \delta.
\]
As $\delta$ is arbitrary, we are done.
\qed
\end{proof}

\begin{definition}\label{def:measurable_set}
* A set $A$ is \underline{measurable} if for each set $E$ we have
\[
m^*E = m^*(A \cap E) + m^*(A^c \cap E).
\]
Note: We know for any set $m^*E \leq m^*(A \cap E) + m^*(A^c \cap E)$ holds. Hence $A$ is measurable iff
\[
m^*E \geq m^*(A \cap E) + m^*(A^c \cap E) \quad \text{for any set } E.
\]
\end{definition}

% --- Page 21 ---
\begin{remark}\label{rem:complement_measurable}
If $A$ is measurable then $A^c$ is measurable.
\end{remark}

\begin{lemma}\label{lem:union_measurable}
If $A_1$ and $A_2$ are measurable sets then $A_1 \cup A_2$ is measurable.
\end{lemma}

\begin{proof}
For any set $E \subseteq X$,

\[
\mu^*(E \cap (A_1 \cup A_2)) = \mu^*(E \cap A_1) + \mu^*(E \cap A_2 \cap A_1^c)
\]

and

\[
\mu^*(E \cap (A_1 \cup A_2)) + \mu^*(E \cap (A_1 \cup A_2)^c) = \mu^*(E \cap A_1) + \mu^*(E \cap A_2 \cap A_1^c) + \mu^*(E \cap A_2^c \cap A_1^c).
\]

Now apply the fact that $A_2$ is a measurable set,

\[
\mu^*(E \cap (A_1 \cup A_2)) + \mu^*(E \cap (A_1 \cup A_2)^c) = \mu^*(E \cap A_1) + \mu^*(E \cap A_1^c) = \mu^*(E)
\]

as $A_1$ is a measurable set. Hence $A_1 \cup A_2$ is measurable.
\end{proof}

% --- Page 22 ---
\begin{lemma}
Let $E$ be any set, and $A_1, \dots, A_n$ a finite sequence of disjoint measurable sets. Then
\[
\mu^*\left(E \cap \left(\bigcup_{i=1}^n A_i\right)\right) = \sum_{i=1}^n \mu^*(E \cap A_i).
\]

\begin{proof}
Base case holds for $n=1$. Assume the given holds for $n-1$ sets.

Observe
\[
E \cap \left(\bigcup_{i=1}^{n-1} A_i\right) \cap A_n = E \cap A_n
\]
and
\[
E \cap \left(\bigcup_{i=1}^{n-1} A_i\right) \cap A_n^c = E \cap \bigcup_{i=1}^{n-1} A_i
\]
as $A_n$ is disjoint from $\left\{A_i\right\}_{i=1}^{n-1}$. Hence
\[
\mu^*\left(E \cap \left(\bigcup_{i=1}^n A_i\right)\right) = \mu^*\left(E \cap \bigcup_{i=1}^{n-1} A_i \cap A_n\right) + \mu^*\left(E \cap \bigcup_{i=1}^{n-1} A_i \cap A_n^c\right)
\]
by measurability of $A_n$. This implies
\[
\mu^*\left(E \cap \left(\bigcup_{i=1}^n A_i\right)\right) = \mu^*(E \cap A_n) + \mu^*\left(E \cap \bigcup_{i=1}^{n-1} A_i\right).
\]
\end{proof}
\end{lemma}

% --- Page 23 ---
By induction hypothesis, we get
\[
\mu^{*}\left(E \cap \left(\bigcup_{i=1}^n A_i\right)\right) = \sum_{i=1}^n \mu^{*}(E \cap A_i).
\]

\begin{theorem}[Carathéodory's Theorem]
If $\mu^*$ is an outer measure on $X$, the collection $\mathcal{M}$ of $\mu^*$-measurable sets is a $\sigma$-algebra and the restriction of $\mu^*$ to $\mathcal{M}$ is a complete measure.
\end{theorem}

\begin{proof}
According to a previous remark and lemma, $\mathcal{M}$ is an algebra and finitely additive. As per the previous lemma,
\[
\mu^{*}\left(E \cap \left(\bigcup_{i=1}^\infty A_i\right)\right) = \sum_{i=1}^\infty \mu^{*}(E \cap A_i)
\]
holds for any $E \subset X$ and collection of $\mu^*$-measurable disjoint sets $\{A_i\}_{i=1}^\infty$.
\end{proof}

% --- Page 24 ---
To ensure the content is properly structured and formatted according to the rules, 
---

In order to show that $\mathcal{M}$ is closed under disjoint unions, let us proceed as follows:

Let $B_n = \bigcup_{j=1}^n A_j$ and $B = \bigcup_{j=1}^\infty A_j$. Hence,
\[
\mu^*(E) = \mu^*(E \cap B_n) + \mu^*(E \cap B_n^c)
\]
\[
\geq \sum_{j=1}^n \mu^*(E \cap A_j) + \mu^*(E \cap B^c).
\]

Observe the statement above holds for any arbitrary $n$. Therefore,
\[
\mu^*(E) \geq \sum_{j=1}^\infty \mu^*(E \cap A_j) + \mu^*(E \cap B^c)
\]
\[
\geq \mu^*\left(E \cap \left(\bigcup_{j=1}^\infty A_j\right)\right) + \mu^*(E \cap B^c)
\]
\[
= \mu^*(E \cap B) + \mu^*(E \cap B^c) \geq \mu^*(E).
\]

It follows that $B \in \mathcal{M}$ and $\mu^*(B) = \sum_{j=1}^\infty \mu^*(A_j)$, so $\mu^*$ is countably additive, by substituting $E = B$.

Finally, if $\mu^*(A) < 0$ for any $E \subset X$, we have
\[
\mu^*(E) \leq \mu^*(E \cap A) + \mu^*(E \cap A^c)
\]
\[
= \mu^*(E \cap A^c) \leq \mu(E).
\]

---

% --- Page 25 ---
---

So, $A \in \mathcal{M}$. Hence $\mu^*|_{\mathcal{M}}$ is complete measure.

\begin{definition}\label{def:premeasure}
Let $A \subset \mathcal{P}(X)$ be an algebra, a function
$\mu_0 \colon A \to [0, \infty]$ will be called a \underline{premeasure} if
\begin{enumerate}
    \item $\mu_0(\emptyset) = 0$,
    \item if $\{A_j\}_{j \in \mathbb{N}}$ is a sequence of disjoint sets in $A$, then
    \[
    \mu_0\left(\bigcup_{j=1}^\infty A_j\right) = \sum_{j=1}^\infty \mu_0(A_j).
    \]
\end{definition}

\begin{remark}\label{rem:measure_vs_premeasure}
Note: The difference between a measure and a premeasure is the domain.
\begin{itemize}
    \item Measure is defined on a $\sigma$-algebra.
    \item Premeasure is defined on an algebra.
\end{itemize}
\end{remark}

\begin{remark}\label{rem:finite_premeasure}
We have a notion of finite and $\sigma$-finite premeasures as same as that of a measure.
\end{remark}

---

% --- Page 26 ---
\begin{remark}
If $\mu_0$ is a premeasure on $\mathcal{A} \subseteq \mathcal{P}(X)$, it induces an outer measure on $X$ in accordance with Proposition~\ref{prop:outer_measure_construction}, namely
\[
\mu^*(E) = \inf \left\{ \sum_{j=1}^\infty \mu_0(A_j) \mid A_j \in \mathcal{A}, \, E \subseteq \bigcup_{j=1}^\infty A_j \right\}.
\end{remark}

\begin{lemma}
If $\mu_0$ is a premeasure on $\mathcal{A}$ and $\mu^*$ as defined in the remark above, then
\begin{enumerate}
    \item $\mu^*|_{\mathcal{A}} = \mu_0$;
    \item every set in $\mathcal{A}$ is $\mu^*$-measurable.
\end{enumerate}
\end{lemma}

\begin{proof}
\begin{enumerate}
    \item[(a)] Suppose $E \in \mathcal{A}$. If $E \subseteq \bigcup_{j=1}^\infty A_j$ with any countable collection $\{A_j\}_{j=1}^\infty \subseteq \mathcal{A}$, then we let $B_n = E \cap \left(A_n \setminus \bigcup_{j=1}^{n-1} A_j\right)$. Observe $B_n$ is a disjoint member of $\mathcal{A}$ whose union is $E$, so
\end{proof}

% --- Page 27 ---
---
\[
\mu_0(E) = \sum_{n=1}^\infty \mu_0(B_n) \leq \sum_{n=1}^\infty \mu_0(A_j).
\]
Since this holds for any arbitrary cover of \(E\), hence as per the definition of \(\mu^*\), we get
\[
\mu_0(E) \leq \mu^*(E).
\]
Now as \(E \in \mathcal{A}\), we can choose a cover \(E \subset \bigcup_j A_j\) s.t. \(A_1 = E\) and \(A_j = \emptyset\) for \(j > 1\). Hence \(\mu_0(E) = \mu^*(E)\).

\begin{proof}
\begin{enumerate}
\item If \(A \in \mathcal{T}\), \(E \subset X\) and \(\varepsilon > 0\), then there is a sequence \(\{B_j\}_{j \in \mathbb{N}} \subset \mathcal{A}\) with
\[
E \subset \bigcup_{j=1}^\infty B_j \quad \text{and} \quad \sum_{j=1}^\infty \mu_0(B_j) \leq \mu^*(E) + \varepsilon.
\]
Since \(\mu_0\) is additive on \(\mathcal{A}\),
\[
\mu_0\left(\bigcup_{j=1}^\infty B_j \cap X\right) = \mu_0\left(\bigcup_{j=1}^\infty B_j \cap (A \cup A^c)\right)
\]
\[
= \mu_0\left(\bigcup_{j=1}^\infty (B_j \cap A)\right) + \mu_0\left(\bigcup_{j=1}^\infty (B_j \cap A^c)\right)
\]
\[
= \sum_{j=1}^\infty \mu_0(B_j \cap A) + \sum_{j=1}^\infty \mu_0(B_j \cap A^c).
\]
Hence,
\[
\mu^*(E) + \varepsilon \geq \sum_{j=1}^\infty \mu_0(B_j \cap A) + \sum_{j=1}^\infty \mu_0(B_j \cap A^c)
\]
\[
\geq \mu_0(E \cap A) + \mu_0(E \cap A^c).
\]
\end{enumerate}
\end{proof}
---

% --- Page 28 ---
As the set $E$ is arbitrary, $A$ is $\mu^*$-measurable.

\begin{theorem}\label{thm:measure_extension}
Let $A \subseteq \mathcal{P}(X)$ be an algebra, $\mu_0$ is a premeasure defined on $A$. Let $\mathcal{M}$ be $\sigma$-algebra generated by $A$. There exists a measure $\mu$ on $\mathcal{M}$ such that $\mu|_A = \mu_0$.
\end{theorem}

\begin{proof}
If we define $\mu^*$ as defined earlier and $\mu := \mu^*|_{\mathcal{M}}$. Then we know by the previous lemma $\mu^*|_A = \mu|_A = \mu_0$ and all $A$ are $\mu^*$-measurable. This implies $\sigma$-algebra of $\mu^*$-measurable sets include $A$.
\end{proof}

% --- Page 29 ---
\begin{theorem}\label{thm:measure_extension}
If $\nu$ is another measure defined on $\mathfrak{M}$ (as per the previous theorem) that extends on $\mu_0$, then $\nu(E) \leq \mu(E)$ for any $E \in \mathfrak{M}$ and equality when $\mu(E) < +\infty$. If $\mu_0$ is $\sigma$-finite, then $\mu$ is the unique extension of $\mu_0$ to a measure on $\mathfrak{M}$.
\end{theorem}

\begin{proof}
If $E \in \mathfrak{M}$ and $E \subset \bigcup_{j=1}^\infty A_j$ where $A_j \in \mathcal{A}$ for all $j \in \mathbb{N}$, then
\[
\nu(E) \leq \sum_{j=1}^\infty \nu(A_j) = \sum_{j=1}^\infty \mu_0(A_j)
\]
holds true for any $\{A_j\}_{j=1}^\infty$ collection. Hence,
\[
\nu(E) \leq \mu(E).
\]

Let $A = \bigcup_{j=1}^\infty A_j$. Now
\[
\nu(A) = \nu\left(\bigcup_{j=1}^\infty A_j\right) = \sum_{j=1}^\infty \nu(A_j) = \lim_{n \to \infty} \sum_{j=1}^n \nu(A_j) = \lim_{n \to \infty} \nu\left(\bigcup_{j=1}^n A_j\right).
\]
Since $\nu$ extends $\mu_0$,
\[
= \lim_{n \to \infty} \mu_0\left(\bigcup_{j=1}^n A_j\right) = \sum_{j=1}^\infty \mu_0(A_j) = \mu_0(A).
\]
\end{proof}

% --- Page 30 ---
---

If $\mu(E) < +\infty$. For every $\varepsilon > 0$ there exist a collection of sets $\{A_\varepsilon, j\}_{j=1}^\infty$, such that $A_\varepsilon = \bigcup_{j=1}^\infty A_{\varepsilon j}$ and $\mu(A_\varepsilon) < \mu^*(E) + \varepsilon$, hence $\mu(A_\varepsilon \setminus E) < \varepsilon$, and $\mu^*(CE) \leq \mu(A) = \nu(A) = \nu(E) + \nu(A \setminus E) \leq \nu(CE) + \mu(A \setminus E) \leq \nu(CE) + \varepsilon$ holds for all $\varepsilon > 0$. This implies $\mu^*(CE) \leq \nu(CE)$.

\begin{proof}
Finally suppose $X = \bigcup_{j=1}^\infty A_j$ with $\mu_0(A_j) < \infty$ where we can assume $A_j$'s are disjoint (or reconstruct a given collection into that way). Then for any $E \in \mathcal{M}$

\[
\mu(E) = \sum_{j=1}^\infty \mu(E \cap A_j) = \sum_{j=1}^\infty \nu^*(E \cap A_j) = \nu^*(CE),
\]

so $\mu = \nu$.
\end{proof}

---

% --- Page 31 ---
% finite disjoint unions of h-intervals is an algebra.

% Borel Measure on $\mathbb{IR}$

\begin{definition}\label{def:h-intervals}
Let $(a,b]$ be denoted as $h$-intervals.
\end{definition}

Let $\mu$ be the Borel measure on $\mathcal{B}_{\mathbb{IR}}$. We define the distribution function as
\[ F(x) = \mu((-\infty, x]). \]

\begin{proposition}\label{prop:premeasure}
Let $F: \mathbb{IR} \to \mathbb{IR}$ be an increasing and right continuous function. If $\{(a_j, b_j]\}_{j=1}^n$ are disjoint $h$-intervals, let
\[ \mu_0\left(\bigcup_{j=1}^n (a_j, b_j]\right) = \sum_{j=1}^n \mu_0((a_j, b_j]), \]
and let $\mu_0(\emptyset) = 0$. Then $\mu_0$ is a premeasure on algebra $\mathcal{A}$.
\end{proposition}

\begin{proof}
(a) $\mu_0$ is well defined.

Given any collection of disjoint $h$-intervals, if $\{(a_j, b_j]\}_{j=1}^n$ by relabelling, we get
\[ a = a_1 < b_1 = a_2 < \cdots < b_n = b. \]
This implies
\[ \sum_{j=1}^n [F(b_j) - F(a_j)] = F(b) - F(a). \]

% --- Page 32 ---
More generally, if $\{I_i\}_{i=1}^n$ and $\{J_j\}_{j=1}^m$ are finite sequences of disjoint $h$-intervals such that $\bigcup_{i=1}^n I_i = \bigcup_{j=1}^m J_j$, by the reasoning above we can say,

\[
\sum_{i} \mu_0(I_i) = \sum_{i,j} \mu_0(I_i \cap J_j) = \sum_{j} \mu_0(J_j).
\]

\begin{proof}
(b) $\mu_0$ is countably additive

Observe $\mu_0$ is finitely additive by construction. We want to show if $\sum I_i \in \mathcal{U}$ s.t. $\bigcup_{i \in \mathbb{N}} I_i \in \mathcal{U}$ then

\[
\mu_0\left(\bigcup_{i=1}^\infty I_i\right) = \sum_{i=1}^\infty \mu_0(I_i).
\]

As $\bigcup_{i=1}^\infty I_i \in \mathcal{U}$, we can separate $\{I_i\}_{i=1}^\infty$ into finitely many subsequences whose union is an $h$-interval. Assume $\bigcup_{i=1}^n I_i$ is an $h$-interval (otherwise you look into a smaller collection of $I_i$'s whose union is an $h$-interval.)

Let $I = \bigcup_{i=1}^\infty I_i = [a, b)$. Observe

\[
\mu_0(I) = \mu_0\left(\bigcup_{j=1}^n I_j\right) + \mu_0\left(I \setminus \bigcup_{j=1}^n I_j\right)
\]

\[
\geq \mu_0\left(\bigcup_{j=1}^n I_j\right) = \sum_{j=1}^n \mu_0(I_j).
\]
\end{proof}

% --- Page 33 ---
To show the reverse inequality, consider the following case:

\begin{proof}
\textbf{Case (ci):} If $a$, $b$ are finite.

As $F$ is right continuous, for any $\varepsilon > 0$, there exists $\delta_\varepsilon > 0$ such that
\[ F(a + \delta_\varepsilon) - F(a) < \varepsilon. \]

If $I_j = (a_j, b_j]$ for $j \in \mathbb{N}$, then we can find a $\delta_{\varepsilon, j} > 0$ such that
\[ F(b_j + \delta_{\varepsilon, j}) - F(b_j) < \frac{\varepsilon}{2^j}. \]

Now, $\left\{ (a_j, b_j + \delta_{\varepsilon, j}) \right\}_{j=1}^\infty$ is a collection of open intervals covering $[a + \delta_\varepsilon, b]$, so there is a finite subcover. Let $(a_1, b_1 + \delta_1), \ldots, (a_N, b_N + \delta_N)$ cover $[a + \delta_\varepsilon, b]$ and $b_i + \delta_i \in (a_{i+1}, b_{i+1} + \delta_{i+1})$. (We have replaced $\delta_{\varepsilon, i} = \delta_i$ for all $i = \{1, \ldots, N-1\}$.)

For any $\varepsilon > 0$, observe
\[
\mu_0(I) < F(b) - F(a + \delta_\varepsilon) + \varepsilon
\]
\[
\leq F(b_N + \delta_N) - F(a_1) + \varepsilon.
\]
\end{proof}

% --- Page 34 ---
\[
= F(b_N + \delta_N) - F(a_N) + \sum_{i=1}^{N-1} \left[ F(a_{j+1}) - F(a_j) \right] + \varepsilon
\]

\[
\leq F(b_N + \delta_N) - F(a_N) + \sum_{i=1}^N \left[ F(b_j + \delta_j) - F(a_j) \right] + \varepsilon
\]

\[
< \left( \sum_{i=1}^N F(b_j) + \frac{\varepsilon}{2^j} - F(a_j) \right) + \varepsilon < \sum_{i=1}^\infty \mu(I_j) + 2\varepsilon.
\]

As the above statement holds for all $\varepsilon > 0$, we are done.

\begin{itemize}
    \item[(ii)] If one of them is NOT finite.

    Say $a = -\infty$. Fix $\varepsilon > 0$. Then for any $M > 0$, we can find an open cover $(a_j + b_j + S_j)$ covering $[-M, b]$. Hence by the same reasoning we get

    \[
    F(b) - F(-M) \leq \sum_{j=1}^\infty \mu_0(I_j) + 2\varepsilon,
    \]

    whereas if $b = \infty$, for any $M < \infty$ we get

    \[
    F(M) - F(a) \leq \sum_{j=1}^\infty \mu_0(I_j) + 2\varepsilon.
    \]

    We get the desired result by $\varepsilon \to 0$ followed by $M \to \infty$. (Go through this).
\end{itemize}

% --- Page 35 ---
\begin{theorem}\label{thm:measure_induced_by_increasing_function}
If $F: \mathbb{R} \to \mathbb{R}$ is any increasing, right continuous function, tpremeasure_induction}. It is clear that $F$ and $G$ induce the same premeasure if and only if $F - G$ is constant, and these premeasures are $\sigma$-finite as $\mathbb{R} = \bigcup_{j \in \mathbb{Z}} (j, j+1]$.
\end{proof}
\end{theorem}

% --- Page 36 ---
The first two assertions therefore follow from \cref{thm:slug}.

As for the converse, the monotonicity of $\mu$ implies $F$ to be an increasing function.

Observe $\mu$ is continuous from above and below implies the right continuity of $F$ for $x \geq 0$ and $x < 0$.

It is evident that $\mu = \mu_F$ on $\mathcal{U}$, and hence $\mu = \mu_F$ on $\mathcal{B}_{\mathbb{R}}$ by uniqueness of \cref{thm:slug}.

\begin{remark}\label{rem:theory_development}
    \begin{enumerate}
        \item The theory could have also been developed for $[a,b)$ intervals and $F$ as left continuous.
        \item If $\mu$ is a finite Borel measure on $\mathbb{R}$ then
        \[
        F(x) = \mu((-\infty, x])
        \]
        is the cumulative distribution function.
        \item Any increasing right continuous $F$ gives us a complete measure.
    \end{enumerate}
\end{remark}

% --- Page 37 ---
We denote this complete measure whose domain includes $\mathcal{B}_{\mathbb{R}}$ as Lebesgue--Stieltjes measure associated to $F$.

Given an increasing, right continuous function $F$, we get $(\mathbb{R}, \mathcal{M}_\mu, \mu)$ as a complete measure space where for any $E \in \mathcal{M}_\mu$,

\[
\mu(E) = \inf \left\{ \sum_{j=1}^\infty [F(b_j) - F(a_j)] \mid E \subseteq \bigcup_{j=1}^\infty (a_j, b_j] \right\}.
\]

\begin{lemma}\label{lem:measure_equality}
For any $E \in \mathcal{M}_\mu$,

\[
\mu(E) = \inf \left\{ \sum_{j=1}^\infty \mu((a_j, b_j]) \mid E \subseteq \bigcup_{j=1}^\infty (a_j, b_j] \right\}.
\]

\begin{proof}
Let $\nu \colon \mathcal{M}_\mu \to [0, \infty)$ such that

\[
\nu(E) = \inf \left\{ \sum_{j=1}^\infty \mu((a_j, b_j]) \mid E \subseteq \bigcup_{j=1}^\infty (a_j, b_j] \right\}
\]

and $\mu$ be the measure defined in the above given definition.

Goal: $\nu(E) = \mu(E)$ for all $E \in \mathcal{M}_\mu$.
\end{proof}
\end{lemma}

% --- Page 38 ---
---

\begin{enumerate}
    \item We want to show $\mu(E) \leq \nu(E)$ for any $E \in \mathcal{M}_\mu$.

    Suppose $E \subseteq \bigcup_{j=1}^\infty (a_j, b_j)$ and for each $(a_j, b_j)$ we get $\{I_j^k\}_{k=1}^\infty$ as disjoint $h$-intervals covering $(a_j, b_j)$ for each $k$. Let $I_j^k = [c_j^k, c_{j}^{k+1}]$ where $c_j^1 = a_j$ and $c_j^k \to b_j$ as $k \to \infty$. Thus $E \subseteq \bigcup_{j,k} I_{j,k}$. Hence
    \[
    \sum_{j=1}^\infty \mu((a_j, b_j)) = \sum_{j=1}^\infty \mu(I_j^k) \geq \mu(E).
    \]
    (Or you can use $E \subseteq \bigcup_{j \in \mathbb{N}} (a_j, b_j)$ and by monotonicity and subadditivity and apply infimum, we get $\mu(E) \leq \nu(E)$.)

    \begin{proof}
    Since $E \subseteq \bigcup_{j,k} I_{j,k}$ and $\{I_{j,k}\}$ are disjoint, by the subadditivity of $\mu$ we have
    \[
    \mu(E) \leq \sum_{j,k} \mu(I_{j,k}).
    \]
    For each $j$, since $\{I_j^k\}_{k=1}^\infty$ covers $(a_j, b_j)$ and $\mu$ is monotone, we have
    \[
    \mu((a_j, b_j)) \leq \sum_{k=1}^\infty \mu(I_j^k).
    \]
    Summing over $j$ gives
    \[
    \sum_{j=1}^\infty \mu((a_j, b_j)) \leq \sum_{j,k} \mu(I_{j,k}).
    \]
    Since $\mu(E) \leq \sum_{j,k} \mu(I_{j,k})$ and $\sum_{j=1}^\infty \mu((a_j, b_j))$ is the infimum over all such coverings, it follows that $\mu(E) \leq \nu(E)$.
    \end{proof}

    \item We want to show $\nu(E) \leq \mu(E)$; for any $\varepsilon > 0$, we get $\{ (a_j, b_j) \}_{j=1}^\infty$ covering $E$ and
    \[
    \sum_{j=1}^\infty \mu((a_j, b_j)) \leq \mu(E) + \varepsilon.
    \]
\end{enumerate}

% --- Page 39 ---
---

As $F$ is right continuous and increasing, for each $\varepsilon > 0$, we find $\delta_{\varepsilon j} = \delta_j$ for each $j \in \mathbb{N}$ such that
\[ F(b_j + \delta_j) - F(b_j) < \frac{\varepsilon}{2^j}. \]

Then $E \subset \bigcup_{j=1}^\infty (a_j + b_j + \delta_j)$ and
\[
\mu((a_j, b_j + \delta_j)) \leq \mu((a_j, b_j + \delta_j)) = F(b_j + \delta_j) - F(a_j)
\]
\[
= F(b_j + \delta_j) - F(b_j) + F(b_j) - F(a_j)
\]
for all $j \in \mathbb{N}$. Now by taking a sum,
\[
\sum_{j=1}^\infty \mu((a_j, b_j)) \leq \sum_{j=1}^\infty \left[ F(b_j + \delta_j) - F(b_j) \right] + \sum_{j=1}^\infty \mu((a_j, b_j))
\]
\[
\leq \varepsilon + \varepsilon + \mu(E).
\]

\begin{proof}
Observe that $\nu(E) \leq \sum_{j=1}^\infty \mu((a_j, b_j)) \leq \mu(E) + 2\varepsilon$ holds for all $\varepsilon > 0$. Hence $\nu(E) \leq \mu(E)$.
\end{proof}

% --- Page 40 ---
\begin{theorem}\label{thm:measure_characterization}
If $E \in \mathcal{M}_\mu$, then
\[
\mu(E) = \inf \left\{ \mu(U) \mid U \supseteq E \text{ and } U \text{ is open} \right\}
\]
\[
= \sup \left\{ \mu(K) \mid K \subseteq E \text{ and } K \text{ is compact} \right\}.
\end{theorem}

\end{document}